\section{Introduction}
A quadrotor can pull a cable-suspended payload but cannot push it. If all quadrotors are between a moving payload and a wall, changing the tensions cannot brake the payload; a quadrotor first has to swing to its far side (Fig.~\ref{fig:scene}). This delay is hidden by a high-order control barrier function (HOCBF) that tests only the currently achievable acceleration. We quantify the delay and construct a safety certificate from the maximum wallward excursion of a feasible braking maneuver.

Input-constrained HOCBFs~\cite{xiao2022high,breeden2021high}, feasibility constraints~\cite{xiao2022sufficient}, and backup, constructive, and predictive CBFs~\cite{chen2021backup,vanwijk2025,squires2018constructive,breeden2022predictive} address feasibility through different safe sets; backup CBFs keep racing drones inside a geofence~\cite{singletary2022onboard}, and the braking-distance barrier of adaptive cruise control~\cite{ames2014control} is the special case of a constant authority. Cable-robot wrench calculations describe forces available at a fixed geometry, as intervals or zonotopes~\cite{gouttefarde2011interval,bouchard2010ability,erskine2019wrench}; aerial transport models describe how that geometry moves~\cite{sreenath2013dynamics,lee2018geometric}, and CBFs have been used to plan cable transport~\cite{pallar2025optimal}. Concurrently,~\cite{shi2026escape} certifies a quadrotor that must reorient its thrust before it can brake, with a body-rate model of first order; here the authority follows from the cable geometry, the swing is of second order, and the same calculation gives an outer bound. Closed-form backup control is also studied in~\cite{vanwijk2025}; our closed form concerns the future wall excursion for this particular authority model.

\begin{figure}[t]
\centering

\resizebox{\columnwidth}{!}{%
\begin{tikzpicture}[x=1cm,y=1cm,>=stealth,font=\small]

% wall
\fill[pattern=north east lines] (7.6,0.0) rectangle (8.0,4.20);

\draw[very thick] (7.6,0.0) -- (7.6,4.20);

\node[anchor=south] at (7.8,4.20) {wall};

\draw[->] (6.5,2.31) -- (7.5,2.31)
    node[midway,anchor=south] {$n$};

% payload
\filldraw[fill=gray!35,draw=black] (3.05,0.75) rectangle (3.55,1.25);

\node[anchor=north west,inner sep=2pt] at (3.6,0.78) {payload};

\coordinate (A) at (3.3,1.25);

% vertical through the payload and the angle of the cables
\draw[densely dotted] (A) -- ++(0,2.90);

\draw (A) ++(90:1.0) arc[start angle=90,end angle=60,radius=1.0];

\node at ($(A)+(75:1.5)$) {$30^\circ$};

% cables and quadrotors, leaning toward the wall
\draw[thick] (A) -- ++(60:2.60) coordinate (Q);

\draw[very thick] ($(Q)+(-0.38,0)$) -- ($(Q)+(0.38,0)$);

\draw ($(Q)+(-0.38,0.09)$) ellipse [x radius=0.2,y radius=0.05];

\draw ($(Q)+(0.38,0.09)$) ellipse [x radius=0.2,y radius=0.05];

\node[anchor=west] at ($(Q)+(0.65,0.04)$) {quadrotors};

\draw[->,thick] ($(A)+(0.22,-0.02)$) -- ++(60:1.17)
    node[anchor=west,inner sep=2pt] {$T_iq_i$};

% cables and quadrotors after the swing
\draw[densely dashed,gray] (A) -- ++(120:2.60) coordinate (G);

\draw[densely dashed,gray,very thick] ($(G)+(-0.38,0)$) -- ($(G)+(0.38,0)$);

\draw[gray] ($(G)+(-0.38,0.09)$) ellipse [x radius=0.2,y radius=0.05];

\draw[gray] ($(G)+(0.38,0.09)$) ellipse [x radius=0.2,y radius=0.05];

\draw[dashed,->] ($(A)+(76:2.60)$)
    arc[start angle=76,end angle=104,radius=2.60];

\node[anchor=south] at ($(A)+(90:2.92)$) {swing needed to brake};

% velocity and braking direction
\draw[->,thick] (3.65,1.0) -- (4.7,1.0)
    node[anchor=west,inner sep=2pt] {$v$};

\draw[->,thick] (2.95,1.0) -- (1.9,1.0)
    node[anchor=east,inner sep=2pt] {$y=-n$};

% distance to the wall
\draw[densely dotted] (3.3,0.75) -- (3.3,0.1);

\draw[<->] (3.3,0.25) -- (7.6,0.25)
    node[midway,anchor=north] {$h$};

\end{tikzpicture}%
}

\caption{Two quadrotors carry a payload toward a wall. With both cables leaning toward the wall, positive tensions cannot brake it; a cable must first swing to the far side (dashed).}

\label{fig:scene}

\end{figure}

We let each quadrotor reserve part of its thrust for the cable tension and part for swinging the cable; the braking authority that this guarantees is the support function of the zonotope of tension forces and depends on the cable directions alone (Sec.~\ref{sec:theory}). When every cable swings to its limit, the payload covers at most the stopping distance $\Dres$ toward the wall, which is computed exactly (Sec.~\ref{sec:barrier}). The authority barrier is $H:=h-\Dres$: the states with $H\ge0$ and admissible swing states form a controlled-invariant set, in which the braking maneuver satisfies every constraint of the filter of Sec.~\ref{sec:filter}; two further barriers bound the altitude of the payload during the maneuver. With upper bounds on the authority and on the swing acceleration, the same calculation gives an optimistic stopping distance $\Drel$, below which no state is viable (Theorem~\ref{thm:sandwich}). Both bounds hold for every system whose braking authority is bounded by nondecreasing functions of coordinates that the inputs drive only through their second derivative (Sec.~\ref{sec:class}); the cable system is one instance.

The contributions are: (i) a condition under which an initially feasible HOCBF filter must lose feasibility before a cable reaches the braking half-space (Theorem~\ref{thm:negative}); (ii) an explicit stopping distance for the input-constrained taut-cable model, the authority and altitude barriers built from it, and the certified set on which the barrier filter is always feasible (Theorem~\ref{thm:sandwich}(i), Proposition~\ref{prop:instance}); (iii) an explicit outer bound on the viability kernel (Theorem~\ref{thm:sandwich}(ii)); and (iv) simulations that test the certificate on its own model and measure its conservatism and the failure of a direct sampled implementation. In simulation, the HOCBF filter lost feasibility in all 690 trials; on the exact model with a $1\,\mathrm{ms}$ hold the authority filter reached the wall in one of 100 trials started at the boundary of its certified set, by $2\,\mu\mathrm m$; on the full-order model with a $5\,\mathrm{ms}$ hold and the wall barrier alone it reached the wall in 59 of 100. The sampled results concern the implementation and give no safety guarantee.

